\documentclass[12pt]{article}


% --- math/fonts ---
\usepackage{amsmath,amssymb}
\usepackage{braket}
\usepackage[T1]{fontenc}
\usepackage{libertine}
\usepackage[utf8]{inputenc}

% --- graphics/tikz ---
\usepackage{tikz}
\usepackage{qcircuit}

% --- other stuff ---
\usepackage[margin=1in]{geometry}
\usepackage{pdfpages}
\usepackage{multicol}
\usepackage[shortlabels]{enumitem}
\usepackage{cancel}
\usepackage{graphicx}
\usepackage{float}
\usepackage{listings}
\usepackage{booktabs}
\lstset{
    language=Python,
    basicstyle=\ttfamily\footnotesize,
    keywordstyle=\bfseries,
    frame=single,
    breaklines=true,
    keepspaces=true,
    columns=fullflexible,
    showstringspaces=false
}
\let\epsilon\varepsilon
\usepackage{enumitem}
\setlist[enumerate]{label=\arabic*.}
\usepackage[authoryear,round]{natbib}
\usepackage[hidelinks]{hyperref}


\newcommand{\class}{CSCE883 Machine Learning}           
\newcommand{\examnum}{Homework 3}      
\newcommand{\examdate}{10/11/2026}        
\newcommand{\timelimit}{}             




\begin{document} 
\pagestyle{plain}
\thispagestyle{empty}

\noindent
\begin{tabular*}{\textwidth}{l @{\extracolsep{\fill}} r @{\extracolsep{6pt}} l}
 \textbf{\class} & \textbf{Name:} & \textit{Taylor Belcher}\\       
\textbf{\examnum} &&\\
\textbf{\examdate} &&\\
\end{tabular*}\\
\rule[2ex]{\textwidth}{2pt}
% ---

 
\begin{enumerate}
   \item \emph{(Question 1A.)} We have a linear network with $L$ layers with the initial input $h_0 = x$ and at each subsesquent layer $l$ (where $1 \leq l \leq L$), $h_l = \sigma (w_lh_{l-1})$ and $\sigma(z) = \dfrac{1}{1+e^{-z}}$. $|w_l| \leq 1$ and there are no biases. 
   
   \begin{enumerate}[(i)]
    \item We want to estimate an upper bound on $\left|\dfrac{\partial h_L}{\partial x}\right|$ in terms of $L$. We have that 
    \[
    \begin{aligned}
        h_L &= \sigma(w_L h_{L-1}) \\
            &= \sigma\bigl(w_L \sigma\bigl(w_{L-1} \sigma(\cdots \sigma(w_1 x))\bigr)\bigr).
    \end{aligned}
    \]

    At any given level $l \geq 2$, let $g_l(x)$ denote the remaining nested expression, including the weight $w_l$, all the way down to the input $x$:
    \[
        g_l(x) = w_l \sigma\bigl(w_{l-1} \sigma(\cdots \sigma(w_1 x))\bigr),
        \qquad g_1(x) = w_1 x.
    \]
    Thus $h_l = \sigma(g_l(x))$. By the chain rule, the partial derivative of $h_L$ with respect to $x$ can be written as

    \[
    \begin{aligned}
        \left|\frac{\partial h_L}{\partial x}\right|
            &= \left|\sigma'(g_L(x))\,g_L'(x)\right| \\
            &= \left|\sigma'(g_L(x))\right|\,
               \left|w_L\,\sigma'(g_{L-1}(x))\,g_{L-1}'(x)\right| \\
            &= \left|\sigma'(g_L(x))\right|\,
               \left|w_L\,\sigma'(g_{L-1}(x))\right| \\
            &\quad {}\cdot \left|w_{L-1}\,\sigma'(g_{L-2}(x))\right|
               \cdots
               \left|w_2\,\sigma'(g_1(x))\right|\,|w_1| \\
            &= \left(\prod_{l=1}^{L}|w_l|\right)
               \left(\prod_{l=1}^{L}\left|\sigma'(g_l(x))\right|\right) \\
            &\leq 1\cdot \prod_{l=1}^{L}\left|\sigma'(g_l(x))\right|.
    \end{aligned}
    \]
    The inequality follows from $|w_l| \leq 1$ for every layer, so $\prod_{l=1}^{L}|w_l| \leq 1$. Now the derivative of the sigmoid activation function is
    \[
        \sigma'(z) = -(1+e^{-z})^{-2}(-e^{-z})
        = \frac{e^{-z}}{(1+e^{-z})^2}.
    \]
    To bound this derivative, observe that
    \[
        1-\sigma(z) = 1-\frac{1}{1+e^{-z}}
        = \frac{e^{-z}}{1+e^{-z}}.
    \]
    Thus we can rewrite the derivative as
    \[
        \sigma'(z) = \frac{1}{1+e^{-z}}\cdot\frac{e^{-z}}{1+e^{-z}}
        = \sigma(z)\bigl(1-\sigma(z)\bigr).
    \]
    Since $0<\sigma(z)<1$, this derivative is positive. Completing the square gives
    \[
        |\sigma'(z)| = \sigma(z)\bigl(1-\sigma(z)\bigr)
        = \frac14-\left(\sigma(z)-\frac12\right)^2
        \leq \frac14.
    \]
    Applying this bound to each of the $L$ factors in our product yields
    \[
    \begin{aligned}
        \left|\frac{\partial h_L}{\partial x}\right|
            &\leq \prod_{l=1}^{L}\left|\sigma'(g_l(x))\right| \\
            &\leq \prod_{l=1}^{L}\frac14 \\
            &= \boxed{\left(\frac14\right)^L = 4^{-L}}.
    \end{aligned}
    \]
    Therefore the magnitude of the input derivative is bounded above by $4^{-L}$, which decreases exponentially with the number of layers. Increasing the number of layers will decrease the gradient exponentially, making it so that earlier layers will take much longer to train. 

    \item If the weights were not bounded above by 1, then we could counter-balance the derivative of the sigmoid function with a larger product of weights, but this is not a guarantee as the weights would also increase the input to the sigmoid functions and the larger the input, the flatter the curve on the sigmoid gets, which means the derivative would also be getting smaller as it moves into the saturation region. So increasing the weights can help with a vanishing gradient, but it is not a guaranteed solution. More precisely, the weight bound matters because each layer contributes both a weight and a sigmoid derivative:
    \[
        \left|\frac{\partial h_l}{\partial h_{l-1}}\right|
        = \left|w_l\sigma'(w_lh_{l-1})\right|
        \leq \frac{|w_l|}{4}.
    \]
    With $|w_l|\leq 1$, each factor is at most $1/4$, giving the exponential bound in part (i). Without this weight bound, larger weights could offset small sigmoid derivatives, so small sigmoid derivatives alone do not prove that every deep network must have vanishing gradients. However, larger weights can also move the sigmoid inputs into regions where its derivative is small. Increasing the weights is therefore not a guaranteed solution.

    A vector residual block adds its input back to a learned transformation of that input:
    \[
        \mathbf h_{l+1}=\mathbf h_l+F_l(\mathbf h_l),
        \qquad F_l:\mathbb R^d\to\mathbb R^d.
    \]
    Its Jacobian contains the partial derivative of each output component with respect to each input component. Writing the $i$th output component as $h_{l+1,i}=h_{l,i}+F_{l,i}(\mathbf h_l)$ gives
    \[
    \begin{aligned}
        [J_l]_{ij}
        &=\frac{\partial h_{l+1,i}}{\partial h_{l,j}} \\
        &=\frac{\partial h_{l,i}}{\partial h_{l,j}}
          +\frac{\partial F_{l,i}(\mathbf h_l)}{\partial h_{l,j}} \\
        &=\delta_{ij}+[J_{F_l}(\mathbf h_l)]_{ij},
    \end{aligned}
    \]
    where $\delta_{ij}=1$ if $i=j$ and $0$ otherwise, is Kroenecker's delta. These entries form the identity matrix, so
    \[
        \boxed{J_l=I+J_{F_l}(\mathbf h_l)}.
    \]
    For a scalar loss $\mathcal L$, the chain rule propagates its gradient backward through the block as
    \[
        \nabla_{\mathbf h_l}\mathcal L
        =\bigl(I+J_{F_l}(\mathbf h_l)\bigr)^{\mathsf T}
         \nabla_{\mathbf h_{l+1}}\mathcal L.
    \]
    The identity term provides a direct contribution from the next layer's gradient. When $J_{F_l}$ is small, the block's Jacobian is close to $I$, which helps preserve gradients. This does not guarantee successful training: for example, $J_{F_l}=-I$ would cancel the identity term entirely, and products of Jacobians across many blocks can still shrink or amplify gradients.
   \end{enumerate} 

   \item \emph{(Question 1B.)}

   The input is $3\times24\times24$ in channel--height--width order.
   Convolutions have stride 1, no padding, dilation 1, and one bias per
   output channel. Pooling uses a $2\times2$ window with stride 2,
   no padding, and floor rounding.

   \begin{enumerate}[(i)]
    \item For either spatial dimension $n$, the output-size formula is
    \[
        n_{\mathrm{out}}
        =\left\lfloor\frac{n_{\mathrm{in}}+2p-d(k-1)-1}{s}\right\rfloor+1,
    \]
    where $k$ is the kernel size, $p$ is padding, $d$ is dilation,
    and $s$ is stride. Here $p=0$ and $d=1$, so this becomes
    $n_{\mathrm{out}}=\lfloor(n_{\mathrm{in}}-k)/s\rfloor+1$.
    Since height and width are equal, the same calculation applies to both:
    \begin{center}
    \begin{tabular}{l c c}
        \toprule
        Layer & Spatial calculation & Output dimension \\
        \midrule
        Conv 1 & $\lfloor(24-3)/1\rfloor+1=22$ & $12\times22\times22$ \\
        Pool 1 & $\lfloor(22-2)/2\rfloor+1=11$ & $12\times11\times11$ \\
        Conv 2 & $\lfloor(11-3)/1\rfloor+1=9$ & $128\times9\times9$ \\
        Pool 2 & $\lfloor(9-2)/2\rfloor+1=4$ & $128\times4\times4$ \\
        \bottomrule
    \end{tabular}
    \end{center}
    A convolution's channel count is its number of output filters, while
    pooling preserves the channel count. In Pool 2, floor rounding gives
    $4$, so the last row and column of the $9\times9$ input are not covered
    by a pooling window.

    \item Each convolutional output filter spans all input channels.
    Therefore a $k\times k$ filter has $C_{\mathrm{in}}k^2$ weights and
    one bias, giving
    \[
        N_{\mathrm{conv}}=C_{\mathrm{out}}\cdot(C_{\mathrm{in}}k^2+1).
    \]
    For Conv 1, each output filter has $3\cdot3\cdot3=27$ weights
    and one bias. With 12 output filters, the total is
    \[
        N_{\mathrm{Conv\,1}}=12(27+1)=\boxed{336}
        \quad\text{($324$ weights and $12$ biases).}
    \]
    Conv 2 receives 12 channels, so each output filter has
    $12\cdot3\cdot3=108$ weights and one bias. With 128 output filters,
    \[
        N_{\mathrm{Conv\,2}}=128(108+1)=\boxed{13{,}952}
        \quad\text{($13{,}824$ weights and $128$ biases).}
    \]
    Max pooling only selects the maximum within each window and has
    no trainable weights or biases. Thus Pool 1 and Pool 2 each have
    $\boxed{0}$ trainable parameters.

    \item After Pool 2, flattening the $128\times4\times4$ tensor gives
    \[
        n_{\mathrm{flat}}=128\cdot4\cdot4=\boxed{2{,}048}
    \]
    features. A fully connected layer with 10 outputs has one weight
    for each input--output pair and one bias per output, so it adds
    \[
        N_{\mathrm{FC}}=2{,}048\cdot10+10=\boxed{20{,}490}
    \]
    trainable parameters. Flattening itself has no trainable parameters.

    In a convolution, the same filter weights and bias are reused at
    every spatial location. Thus its parameter count depends on the
    kernel size and channel counts, rather than the number of output
    locations. A fully connected mapping learns separate weights for
    every input--output pair. For example, a fully connected mapping
    from the original $3\cdot24\cdot24=1{,}728$ inputs to the
    $12\cdot22\cdot22=5{,}808$ outputs of Conv 1 would require
    \[
        1{,}728\cdot5{,}808+5{,}808=10{,}042{,}032
    \]
    parameters, compared with Conv 1's 336. This reduction comes from
    both local connectivity (each output uses a $3\times3$ spatial patch)
    and weight sharing (all locations use the same filters).
   \end{enumerate}

\end{enumerate}

\end{document}

% Quick aligned table for equations and explanations
%
%\begin{center}
%  \begin{tabular}{l l l l}
%    % template line 
%    % LHS & $=$ & RHS & \emph{(comment)}\\
%   \end{tabular}
% \end{center}
